
% --- SLIDE 3: CONTAINERS ---
\begin{frame}[fragile]{Containers, in One Slide}

    \begin{columns}[T]
        \begin{column}{0.5\textwidth}
            \begin{lstlisting}[basicstyle=\ttfamily\scriptsize]
FROM python:3.12-slim

WORKDIR /app
COPY requirements.lock .
RUN pip install --no-cache-dir \
        --require-hashes -r requirements.lock

COPY serve.py model.skops ./
EXPOSE 8000
CMD ["uvicorn", "serve:app", \
     "--host", "0.0.0.0", "--port", "8000"]
            \end{lstlisting}
        \end{column}
        \begin{column}{0.46\textwidth}
            \small
            \begin{itemize}
                \setlength\itemsep{0.4em}
                \item A container freezes the OS libraries, the Python version and the packages
                      \textbf{together}. That is the part a lockfile alone cannot do.
                \item Refer to images by \textbf{digest} (\texttt{@sha256:\ldots}), not by the tag
                      \texttt{latest}. Tags move; digests do not.
                \item \textbf{Do not bake giant model files into the image} if they change often --- pull them
                      from the registry at start-up, by version.
                \item GPU serving adds a CUDA-matched base image; the driver on the host must be new enough.
            \end{itemize}
        \end{column}
    \end{columns}

    \vspace{0.6em}
    \begin{block}{What a container does \emph{not} give you}
        Reproducible \emph{training}. Same image, unseeded run, different model. Containers pin the
        environment; seeds and data versions pin the result.
    \end{block}

\end{frame}

\subsection{Serving Patterns}
